% concatenated from Bhardwaj_Bhandari.tex
\documentclass[article,12pt]{amsart}
%\documentclass{article}
\usepackage{lineno,hyperref,enumitem}
\usepackage{amsmath,amssymb,dsfont}
%\usepackage[affil-it]{authblk}
%\modulolinenumbers[5]

%\journal{Pacific Journal of Mathematics }

\newtheorem{thm}{Theorem}[section]
\newtheorem{defn}[thm]{Definition}
\newtheorem{prop}[thm]{Proposition}
\newtheorem{cor}[thm]{Corollary}
\newtheorem{lem}[thm]{Lemma}
\newtheorem{rem}[thm]{Remark}
\newtheorem{conj}[thm]{Conjecture}
\newtheorem{ex}[thm]{Example}
\newtheorem{quest}[thm]{Question}
\newtheorem{obs}[thm]{Observation}
\newtheorem{re}[thm]{Remark}
%\newproof{proof}{Proof}

\newcommand{\ds}{\displaystyle}

%%%%%%%%%%%%%%%%%%%%%%%
%% Elsevier bibliography styles
%%%%%%%%%%%%%%%%%%%%%%%
%% To change the style, put a % in front of the second line of the current style and
%% remove the % from the second line of the style you would like to use.
%%%%%%%%%%%%%%%%%%%%%%%

%% Numbered
%\bibliographystyle{model1-num-names}

%% Numbered without titles
%\bibliographystyle{model1a-num-names}

%% Harvard
%\bibliographystyle{model2-names.bst}\biboptions{authoryear}

%% Vancouver numbered
%\usepackage{numcompress}\bibliographystyle{model3-num-names}

%% Vancouver name/year
%\usepackage{numcompress}\bibliographystyle{model4-names}\biboptions{authoryear}

%% APA style
%\bibliographystyle{model5-names}\biboptions{authoryear}

%% AMA style
%\usepackage{numcompress}\bibliographystyle{model6-num-names}

%% `Elsevier LaTeX' style
%\bibliographystyle{elsarticle-num}
%%%%%%%%%%%%%%%%%%%%%%%

\begin{document}



\title[Study of weaving frames in Krein spaces]{Study of weaving frames in Krein spaces}

%----------Author 1
%\author[D. Haldar]{Debasis Haldar}
%\address{Department of Mathematics\\ NIT Rourkela\\ Rourkela 769008\\ India}
%\email{debamath.haldar@gmail.com}
%----------Author 2
%\author[D. Singh]{Divya Singh}
%\address{Department of Mathematics\\ NIT Rourkela\\ Rourkela 769008\\ India}
%\email{divya\_allduniv@yahoo.com}

%----------Author 1


\author[A. Bhardwaj]{Avinash Bhardwaj}


\author[A. Bhandari]{Animesh Bhandari$^*$}
\address{Department of Mathematics\\ SRM University, {\it AP} - Andhra Pradesh\\ Neeru Konda, Amaravati, Andhra Pradesh - 522502\\ India}

\email{bhandari.animesh@gmail.com, animesh.b@srmap.edu.in}
%\email{ avinash_bhardwaj@srmap.edu.in }
%\thanks{Third author was supported by DST-SERB project MTR/2017/000797}
%----------classification, keywords, date
\subjclass[2010]{42C15, 46E22, 46C07, 47A05}

\keywords{frame, weaving frames, fundamental symmetry, Krein space\\ $*$ is the corresponding author}

\begin{abstract}
Inspired by the work of Bemrose et al. \cite{Be16}, we delve into the study of weaving frames in Krein spaces. This paper presents a comprehensive exploration of various properties and characterizations of Krein space weaving frames. In support of our findings, several examples and counter examples are provided, illustrating the applicability of the theoretical results. Additionally, we extend the discussion to an important application in probabilistic erasure, highlighting how weaving frames can be used to mitigate data loss in such scenarios. This work contributes to the broader understanding of frame theory in indefinite inner product spaces.

%Furthermore, the scrutiny substantiates its findings through comprehensive discussions of examples, reinforcing the validity and applicability of the results in the context of frame theory and model spaces.
\end{abstract}



\maketitle

%\linenumbers

%\tableofcontents

\section{Introduction}
The concept of frames, first introduced by Duffin et al. \cite{duffin1952class}, gained widespread attention when Daubechies et al. \cite{2} underscored its importance in signal processing. This renewed interest led to frame theory becoming a crucial tool in various fields, including Mathematics, Physics, and Engineering. Its versatility has since made frames indispensable in different areas, such as, wavelet analysis, quantum mechanics, and communication systems.

We first discuss the basics of frames to establish the foundational concepts required to our results.
 
 
\begin{defn}
A sequence  $\{ f_n \}_{n=1}^\infty$ in a Hilbert space $\mathcal{H}$ is called a \emph{frame} for $\mathcal H$ if there exist constants $A,B >0$ so that 
\begin{equation}\label{Eq:Frame} A\|f\|^2~ \leq ~\sum\limits_{n=1}^\infty |\langle f,f_n\rangle|^2 ~\leq ~B\|f\|^2,
\end{equation} for all $f \in \mathcal{H}$. The numbers $A, B$ are called \emph{frame bounds}. For a given  frame $\{f_n\}_{n=1}^\infty$ for $\mathcal{H}$, the \emph{pre-frame operator} or \emph{synthesis operator} is a bounded linear operator $T: \ell^2(\mathbb N)\rightarrow \mathcal{H}$ and is defined as $T\{c_n\}_{n=1}^\infty = \sum\limits_{n=1}^\infty c_n f_n$. The adjoint of $T$, $T^*: \mathcal{H} \rightarrow \ell^2(\mathbb N)$ is  given by $T^*f = \{\langle f, f_n\rangle\}_{n=1}^\infty$, is called the \emph{analysis operator}. The \emph{frame operator}, $S=TT^*: \mathcal{H}\rightarrow \mathcal{H}$, is defined as 
\begin{equation}
Sf=TT^*f = \sum\limits_{n=1}^\infty \langle f, f_n\rangle f_n~~\emph{for~all~} f\in {\mathcal H}.
\end{equation}
We notice that depending upon the nature of frame bounds we get special kinds of frames, e.g. in the inequality (\ref{Eq:Frame}) if $A=B$ then the given frame is called a tight frame, in addition to the same if $A=B=1$ it is called a Parseval frame. It is well-known that the frame operator is bounded, positive, self adjoint and invertible.

Using the invertibility property, every element in the Hilbert space can be reconstructed as,
\begin{equation}
f=\sum\limits_{n=1}^\infty \langle f, S^{-1}f_n \rangle f_n = \sum\limits_{n=1}^\infty \langle f, f_n \rangle S^{-1}f_n.
\end{equation}

Furthermore, $T^*T\{c_n\}=\sum\limits_{m=1}^\infty \langle f_n, f_m\rangle c_n$ for every $\{c_n\} \in \ell^2(\mathbb N)$ and its matrix representation is $(T^*T)_{m,n} = \langle \delta_m, T^*T \delta_n \rangle_{l^2}$. Thus we have,
\begin{eqnarray*}
(T^*T)_{m,n} = \langle \delta_m, T^*T \delta_n \rangle_{l^2} &=& \langle T \delta_m, T \delta_n \rangle_{\mathcal H}
\\&=& \langle f_m, f_n \rangle_{\mathcal H}
\\&=& G_{m, n}
\end{eqnarray*}
which is known as Gramian matrix. It is to be noted that the Gramian of a frame is invertible if and only if the given frame is a Riesz basis.

For detailed discussion regarding frames, their generalization and applications we refer \cite{christensen2003introduction, 1, Kaf09, Gu18, 3, Po20, Ba18, Fi97, Fi98, Fi06, Fi00, Na12, Na13, Bo06, Cz02, Cz06, Ha07, Ha00, He04, An79}.
\end{defn}

If $(\mathcal K, [.,.])$ is a Krein space with fundamental decomposition $\mathcal K_{\delta^+} \oplus \mathcal K_{\delta^-}$ and  fundamental  symmetry  $J = P_{\delta^+} - P_{\delta^-}$, where $P_{\delta^{\pm}}$ is the orthogonal  projection onto $\mathcal{K}_{\delta^{\pm}}$, then for every $k \in \mathcal K$ we have, $$J(k) = J(k_{\delta^+}+ k_{\delta^-}) = k_{\delta^+} - k_{\delta^-}, \quad k = k_{\delta^+}+ k_{\delta^-}\in \mathcal K_{\delta^+} \oplus \mathcal K_{\delta^-}. $$ In this context the associated $J$-inner product is given by,
$$[k, k']_J = [k, Jk'],  \quad k, k' \in \mathcal{K},$$ and hence ($\mathcal{K}, [., .]_J$) is a Hilbert space. The inner product $[., .],$ which is positive definite, induces a topology on Krein space $\mathcal{K}$ through $J$-norm define as, $$\|k\|_{J}^2 = [k, k]_J = [k, Jk], \quad \forall k \in \mathcal{K}.$$
Here ($\mathcal{K}_{\delta^+}, +[., .]$) is a Hilbert space and ($\mathcal{K}_{\delta^-}, -[., .]$)  is known as anti Hilbert space.

A $J$-orthonormalized system is a family of vectors $\{e_n\}_{n=1}^\infty$ satisfying condition $[e_m, e_n]=\delta_{mn}$ for every $m, n \in \mathbb N$. A $J$-orthonormal basis is a $J$-orthonormalized system that also serves as a Schauder basis, providing unique and stable expansions for vectors. If $\{e_n\}_{n=1}^\infty$ is a $J$-orthonormal basis in $\mathcal K$, then any element $k \in \mathcal K$ can be represented as $k=\sum\limits_{n=1}^\infty \lambda_n [k, e_n] e_n$, where $\lambda_n=[e_n, e_n]=\pm 1$.



A sequence $\{z_n\}_{n=1}^\infty$ in a Krein space $(\mathcal K, [.,.])$ is said to be a frame for $\mathcal K$, if for every $k \in \mathcal K$ there exist positive constants $\alpha, \beta$ we have,
\begin{equation}
\alpha \|k\|_{J}^2 \leq \sum\limits_{n=1}^\infty |[k, z_n]|^2 \leq \beta \|k\|_{J}^2,
\end{equation}
where $\|\|_{J}$ is the standard norm in $(\mathcal K, \langle .,. \rangle)$. It is to be noted that a sequence $\{z_n\}_{n=1}^\infty$ is a frame for a Krein space $(\mathcal K, [.,.])$ if and only if it is a frame for $(\mathcal K, \langle .,. \rangle)$ (see \cite{Ga}). The key difference between $J$-frames and classical frames for Hilbert spaces is that some frames in Hilbert space $(\mathcal K, \langle .,. \rangle)$  are not necessarily $J$-frames in Krein spaces $(\mathcal K, [.,.])$. For detailed discussion regarding the same we refer (Example 3.3, \cite{Gi12}).\\


%We will discuss the Moore-Penrose pseudoinverse, which will help us to analyze and interpret our results.

%\begin{lem}\cite{Be03}\label{Moore} Let $\mathcal H$ and $\mathcal K$ be two Hilbert spaces and $T \in \mathcal L(\mathcal H,\mathcal K)$ be a closed range operator, then the following results hold:
	%\begin{enumerate}
		%\item $TT^{\dag}=P_{R(T)}$, $T^{\dag}T=P_{R(T^*)}$
	%	\item $\frac {\|f\|} {\|T^{\dag}\|} \leq \|T^*f\|$ for all $f \in R(T)$.
		%\item $TT^{\dag}T=T$, $T^{\dag}TT^{\dag}=T^{\dag}$, $(TT^{\dag})^*=TT^{\dag}$, $(T^{\dag}T)^*=T^{\dag}T$.
	%\end{enumerate}
%\end{lem}


%The utilization of the Douglas factorization theorem serves a vital role of facilitating the proof of key results and enhancing the overall robustness of our theoretical framework.

%\begin{thm}(Douglas' factorization theorem \cite{Do66})\label{Thm-Douglas} Let $\mathcal{H}_1, \mathcal{H}_2,$ and $\mathcal H$ be Hilbert spaces and $S\in \mathcal L(\mathcal H_1, \mathcal H)$, $T\in \mathcal L(\mathcal H_2, \mathcal H)$. Then the following statements are equivalent:
%	\begin{enumerate}
	%	\item $R(S) \subseteq R(T)$.
	%	\item $SS^* \leq \alpha TT^*$ for some $\alpha >0$.
		%\item $S = TL$ for some $L \in \mathcal L(\mathcal H_1, \mathcal H_2)$.
%	\end{enumerate}
%\end{thm}

\section{Main Results}

In this section, we explore various properties and characterizations of weaving frames within the structure of Krein spaces. We begin by introducing the concept of Krein space weaving frames, which extends the classical frame theory into indefinite inner product spaces. Our discussion focuses on  the interplay between the positive and negative components of Krein spaces.

\begin{defn}
Let $\{z_n\}_{n=1}^\infty$ and $\{z'_n\}_{n=1}^\infty$ be two frames for the Krein space $(\mathcal K, [.,.])$, then they are said to be weaving frames for $\mathcal K$ if for for every $\sigma \subset \mathbb N$, $\{z_n\}_{n \in \sigma} \cup \{z'_n\}_{n \in \sigma^c}$ is a frame for $\mathcal K$. Thus for every $k \in \mathcal K$ and every $\sigma \subset \mathbb N$, there exist positive constants $\alpha \leq \beta$ we have,
\begin{equation}\label{1}
\alpha \|k\|_{J}^2 \leq \sum\limits_{n \in \sigma} |[k, z_n]|^2 + \sum\limits_{n \in \sigma^c} |[k, z'_n]|^2 \leq \beta \|k\|_{J}^2.
\end{equation}
\end{defn}

We now define the corresponding  frame operator for the associated weaving frames, establishing its roles in the framework.

The weaving frame operator $S:\mathcal K \rightarrow \mathcal K$ is defined as, 
$$S_{\sigma}(k)=TJT^*=\sum\limits_{n \in \sigma} [k, z_n] z_n + \sum\limits_{n \in \sigma^c} [k, z'_n] z'_n,$$ where $T, T^*$ and $J$ are associated synthesis, analysis operators and fundamental symmetry on $(\mathcal K, [.,.])$.

\begin{ex}
let us assume a Krein space $(\mathcal K = \mathbb C^2, [.,.])$ with the associated indefinite inner product defined as $[(k_1 , k_2), (l_1,  l_2)] = k_1\overline l_1 - k_2\overline l_2$, where $\ (k_1, k_2)  \in {\mathcal K}$, $\ (l_1, l_2) \in {\mathcal K}$. Consider two frames 
$$
\{z'_n\}_{n=1}^3\ = \{(1, 1), (0, 1), (1, 0) \}$$ and $$\{z''_n\}_{n= 1}^3\ = \{(-1, -1), (0, -1), (1, 0) \}.
$$ Then it is easy to verify that for every $\sigma \subset \{1, 2, 3\}$, the equation~\eqref{1} is satisfied. Therefore $\{z'_n\}_{n = 1}^3$ and $\{z''_n\}_{n = 1}^3$ are weaving frames for $(\mathcal K = \mathbb C^2, [.,.])$.


On the other hand , if we consider a Krein space  $(\mathcal K = \mathbb C^3, [.,.])$. with the associated indefinite inner product $[. , .]$ defined as 
$$
[ (k_1, k_2, k_3 ), (l_1, l_2, l_3 )] = k_1\overline l_1 + k_2\overline l_2 - k_3\overline l_3,
$$ where $\ (k_1, k_2, k_3) \in {\mathcal K}$, $\ (l_1, l_2, l_3) \in {\mathcal K}$. Let us consider two frames 
$$
\{z'_n\}_{n=1}^3\ = \{(1, 2, 0), (0, \sqrt{2}, 0), (0, 0, 1) \}
$$ and 
$$
\{z''_n\}_{n= 1}^3\ = \{(0, -\sqrt{2}, 0), (-2, 0, 0), (0, 0, 2) \}.
$$
 It is easy to verify that these frames are not weaving frames if we consider $\sigma$ = $\{2\}$, since for $k =(k_1, 0, 0)$ with $k_1 \neq 0$ we have  $\|k\|_{J}^2 =  \|(k_1, 0,0)\|_{J}^2 = |k_1|^2$, where as 
\begin{eqnarray*}
&&\sum\limits_{n \in \sigma} |[k, z'_n]|^2 + \sum\limits_{n \in \sigma^c} |[k, z''_n]|^2 \\&=& |[(k_1, 0, 0), (0, \sqrt{2}, 0)]|^2
+|[(k_1, 0, 0), (0, -\sqrt{2}, 0)]|^2
 + |[(k_1, 0, 0), (0, 0, 2)]|^2 
\\&=& 0.
\end{eqnarray*}
 Thus the left hand inequality of equation (\ref{1}) is not satisfied.
Therefore, the given frames are not weaving frames.


\end{ex}We present both an example and a non-example of Krein space weaving frames to illustrate their defining characteristics and limitations.



Our first result provides a characterization of weaving frames in Krein spaces through the use of the associated fundamental symmetry.

\begin{thm}\label{char}
Let $\{z_n\}_{n=1}^\infty$ and $\{z'_n\}_{n=1}^\infty$ be two frames for the Krein space $(\mathcal K, [.,.])$, then the following are equivalent:

\begin{enumerate}
\item $\{z_n\}_{n=1}^\infty$ and $\{z'_n\}_{n=1}^\infty$ are weaving frames for $\mathcal K$.
\item $\{Jz_n\}_{n=1}^\infty$ and $\{Jz'_n\}_{n=1}^\infty$ are weaving frames for $\mathcal K$.
\item $\{z_n\}_{n=1}^\infty$ and $\{z'_n\}_{n=1}^\infty$ are weaving frames for $(\mathcal K, [.,.]_J)$.
\item $\{Jz_n\}_{n=1}^\infty$ and $\{Jz'_n\}_{n=1}^\infty$ are weaving frames for $(\mathcal K, [.,.]_J)$.
\end{enumerate}
\end{thm}

The validity of the above theorem is demonstrated through the following example.

\begin{ex}
Consider $\mathcal K = \ell^2$ and the indefinite inner product on $\ell^2$ is defined as
$$
[k,k']=\sum\limits_{n=1}^\infty (-1)^n z_n \overline{z'_n},
$$ where $k=\{z_n\}_{n=1}^\infty, k'=\{z'_n\}_{n=1}^\infty \in \ell^2$. Suppose that $K_{\delta^+} = \{\{z_n\}_{n=1}^\infty: z_n=0 \text{~if ~n ~is~ odd}\}$ and $K_{\delta^-} = \{\{z'_n\}_{n=1}^\infty: z_n=0 \text{~if ~n ~is~ even}\}$. Then $\mathcal K = K_{\delta^+} \oplus  K_{\delta^-}$, where $K_{\delta^+}$ and $K_{\delta^-}$ are complete and hence $\mathcal K = \ell^2, [.,.]$ is a Krein space. 

Consider the frames 
$$
\{z_n\}_{n=1}^\infty=\{e_1, 0, e_2, e_2, e_3, e_3, \cdots\}
$$ and 
$$
\{z'_n\}_{n=1}^\infty=\{e_1, e_1, e_2, e_2, e_3, e_3, \cdots\}
$$ for $\ell^2$, where $\{e_n\}$ is the standard orthonormal basis in $\ell^2$. Then it is easy to see that they are weaving frames for $\ell^2$. If 
$$
J=\begin{pmatrix}
-1 & 0 & 0 & 0 & 0 & \cdots\\
0 & 1 & 0 & 0 & 0 & \cdots\\
0 & 0 & -1 & 0 & 0 & \cdots\\
0 & 0 & 0 & 1 & 0 & \cdots\\
 \vdots & \vdots & \vdots & \vdots & \vdots & \ddots
\end{pmatrix},
$$

Then $\{Jz_n\}_{n=1}^\infty = \{-e_1, 0, e_2, e_2, -e_3, -e_3, \cdots\}$ and \\
$\{Jz'_n\}_{n=1}^\infty = \{-e_1, -e_1, e_2, e_2, -e_3, -e_3, \cdots\}$, and therefore it is easy to verify that the remaining conditions of the theorem are satisfied.
\end{ex}

\noindent \underline{\bf Proof of Theorem 2.3:}\\
\noindent  (\underline{1 $\Longleftrightarrow$ 4}) Let $\{z_n\}_{n=1}^\infty$ and $\{z'_n\}_{n=1}^\infty$ be weaving frames for $\mathcal K$ with the universal bounds $\alpha, \beta$. Then for every $k \in \mathcal K$ and every $\sigma \subset \mathbb N$ we have, 
\begin{equation}\label{5}
	\alpha \|k\|_{J}^2 \leq \sum\limits_{n \in \sigma} |[k, z_n]|^2 + \sum\limits_{n \in \sigma^c} |[k, z'_n]|^2 \leq \beta \|k\|_{J}^2.
\end{equation}
Now we prove $J^2=\mathcal I$ and $J=J^*$, where $\mathcal I$ is the identity operator and $J$ is the fundamental symmetry on $\mathcal K$. 

Since $\mathcal K = K_{\delta^+} \oplus  K_{\delta^-}$, then for every $k \in \mathcal K$ can be written as $k=k_{\delta^+} + k_{\delta^-}$ and $J = \pi^+ - \pi^-$, where $ \pi^+$ is the orthogonal projection onto $K_{\delta^+}$ and $ \pi^-$ is the orthogonal projection onto $K_{\delta^-}$. 

Thus for every $k \in \mathcal K$ we have,
\begin{eqnarray*}
	J^2(k)=J(Jk)=J(J(k_{\delta^+} + k_{\delta^-}))&=&J( (\pi^+ - \pi^-)(k_{\delta^+} + k_{\delta^-}))\\
	&=&J(k_{\delta^+} - k_{\delta^-})\\
	&=&  (\pi^+ - \pi^-)(k_{\delta^+} - k_{\delta^-})\\
	&=& k_{\delta^+} + k_{\delta^-})\\
	&=& k.
\end{eqnarray*}
Furthermore, for every $k \in \mathcal K$ we have,
\begin{eqnarray*}
	[J^*(k), k] = [k, J(k)]=[k, (\pi^+ - \pi^-)(k)] &=& [(\pi^+ - \pi^-)^*(k), k]\\
	&=& [(\pi^+ - \pi^-)(k), k]\\
	&=& [J(k), k].
\end{eqnarray*}
Therefore, applying equation (\ref{5}), for every $k \in \mathcal K$ and every $\sigma \subset N$ we have,
\begin{equation*}
	\alpha \|k\|_{J}^2 \leq \sum\limits_{n \in \sigma} |[Jk, Jz_n]|^2 + \sum\limits_{n \in \sigma^c} |[Jk, Jz'_n]|^2 \leq \beta \|k\|_{J}^2.
\end{equation*}
and hence we obtain,
\begin{equation*}
	\alpha \|k\|_{J}^2 \leq \sum\limits_{n \in \sigma} |[k, Jz_n]_{J}|^2 + \sum\limits_{n \in \sigma^c} |[k, Jz'_n]_{J}|^2 \leq \beta \|k\|_{J}^2.
\end{equation*}

\noindent  (\underline{2 $\Longleftrightarrow$ 3}) Let $\{J z_n\}_{n=1}^\infty$ and $\{J z'_n\}_{n=1}^\infty$ be weaving frames for $\mathcal K$ with the universal bounds $\alpha, \beta$. Then for every $k \in \mathcal K$ and every $\sigma \subset \mathbb N$ we have, 
\begin{eqnarray*}
	&& \alpha \|k\|_{J}^2 \leq \sum\limits_{n \in \sigma} |[k, J z_n]|^2 + \sum\limits_{n \in \sigma^c} |[k, J z'_n]|^2 \leq \beta \|k\|_{J}^2.\\
	& \Leftrightarrow & \alpha \|k\|_{J}^2 \leq \sum\limits_{n \in \sigma} |[k, z_n]_{J}|^2 + \sum\limits_{n \in \sigma^c} |[k, z'_n]_{J}|^2 \leq \beta \|k\|_{J}^2.
\end{eqnarray*}
Thus $\{z_n\}_{n=1}^\infty$ and $\{z'_n\}_{n=1}^\infty$ are weaving frames for $(\mathcal K, [.,.]_J)$ and conversely.

\noindent  (\underline{3 $\Longleftrightarrow$ 4}) Let $\{z_n\}_{n=1}^\infty$ and $\{z'_n\}_{n=1}^\infty$ are weaving frames for $(\mathcal K, [.,.]_J)$ with the universal bounds $\alpha, \beta$. 

Then for every $k \in \mathcal K$ and every $\sigma \subset \mathbb N$ we have, 
\begin{equation}\label{6}
	\alpha \|k\|_{J}^2 \leq \sum\limits_{n \in \sigma} |[k, z_n]_{J}|^2 + \sum\limits_{n \in \sigma^c} |[k, z'_n]_{J}|^2 \leq \beta \|k\|_{J}^2.
\end{equation}
For every $k \in \mathcal K$, we prove $\|J(k)\|_{J}=\|k\|_{J}$. Since $\|J\|_{J}=1$, then $\|J(k)\|_{J} \leq \|k\|_{J}$. 

Furthermore, for every $k \in \mathcal K$, $\|k\|_{J} \leq \|J^{-1}\|_{J} \|J(k)\|_{J}$, and hence we have $\|k\|_{J} \leq \|J(k)\|_{J}$.  Thus for every $k \in \mathcal K$, $\|J(k)\|_{J}=\|k\|_{J}$.

Therefore, using equation (\ref{6}) and replacing $k$ by $J(k)$, for every every $\sigma \subset \mathbb N$ we obtain,
\begin{eqnarray*}
	&& \alpha \|J(k)\|_{J}^2 \leq \sum\limits_{n \in \sigma} |[J(k), z_n]_{J}|^2 + \sum\limits_{n \in \sigma^c} |[J(k), z'_n]_{J}|^2 \leq \beta \|J(k)\|_{J}^2\\
	& \Leftrightarrow & \alpha \|k\|_{J}^2 \leq \sum\limits_{n \in \sigma} |[k, J z_n]_{J}|^2 + \sum\limits_{n \in \sigma^c} |[k, Jz'_n]_{J}|^2 \leq \beta \|k\|_{J}^2.
\end{eqnarray*}
This completes the proof.\\


In the following result, we present necessary and sufficient conditions for the existence of Krein space weaving frames. 

\begin{prop}
Let $\{z_n\}_{n=1}^\infty$ and $\{z'_n\}_{n=1}^\infty$ be frames for the Krein space $(\mathcal K, [.,.])$ with the bounds $\alpha_1 \leq \beta_1$ and $\alpha_2 \leq \beta_2$, then the following are equivalent:
\begin{enumerate}
\item $\{z_n\}_{n=1}^\infty$ and $\{z'_n\}_{n=1}^\infty$ are weaving frames for $\mathcal K$.

\item For every $\sigma \subset \mathbb N$, if $S_{\sigma}$ is the corresponding weaving frame operator, then for every $k \in \mathcal K$ there exists $\alpha >0$ so that we have,
$$\|S_{\sigma}(k)\|_{J} \geq \alpha \|k\|_{J}.$$
\end{enumerate}
\end{prop}

\begin{proof}  \noindent  (\underline{1 $\implies$ 2})
	
Let $\{z_n\}_{n=1}^\infty$ and $\{z'_n\}_{n=1}^\infty$ be weaving frames for $(\mathcal K, [.,.])$ with the universal bounds $\alpha, \beta$. Then for every $k \in \mathcal K$ and every $\sigma \subset \mathbb N$ we have,
\begin{equation*}
\alpha \|k\|_{J}^2 \leq \sum\limits_{n \in \sigma} |[k, z_n]|^2 + \sum\limits_{n \in \sigma^c} |[k, z'_n]|^2 \leq \beta \|k\|_{J}^2.
\end{equation*}
and hence applying theorem \ref{char} we obtain,
\begin{equation}\label{2}
\alpha \|k\|_{J}^2 \leq \sum\limits_{n \in \sigma} |[k, z_n]_{J}|^2 + \sum\limits_{n \in \sigma^c} |[k, z'_n]_{J}|^2 \leq \beta \|k\|_{J}^2.
\end{equation}
Thus using equation (\ref{2}), for every $k \in \mathcal K$ and every $\sigma \subset \mathbb N$ we obtain,
\begin{eqnarray*}
\|S_{\sigma}k\|_{J} = \sup\limits_{\|l\|=1} |[S_{\sigma}k, l]_{J}| & \geq & \left [ S_{\sigma}k, \frac{k}{\|k\|_{J}}\right ]_{J}\\
& = & \frac{1} {\|k\|_{J}} [S_{\sigma}k, k]_{J}\\
& \geq & \alpha \|k\|_{J}.
\end{eqnarray*}

 \noindent  (\underline{2 $\implies$ 1}) For every $k \in \mathcal K$ and every $\sigma \subset \mathbb N$, the corresponding weaving frame operator is given by, $S_{\sigma} k = TJT^* k$, where $T, T^*$ and $J$ are associated synthesis, analysis operators and fundamental symmetry on $(\mathcal K, [.,.])$. 
 
 Since $\|J\| = 1$ and $\|S_{\sigma}(k)\|_{J} \geq \alpha \|k\|_{J}$, and $\{z_n\}_{n=1}^\infty$ and $\{z'_n\}_{n=1}^\infty$ are frames for the Krein space $(\mathcal K, [.,.])$ with the bounds $\alpha_1 \leq \beta_1$ and $\alpha_2 \leq \beta_2$ then we have,
 \begin{equation}\label{3}
 \alpha^2 \|k\|_{J}^2 \leq \|T\|_{J}^2 \|T^*k\|_{J}^2 \leq \|T\|_{J}^2 (\beta_1 + \beta_2) \|k\|_{J}^2.
\end{equation}  
Again we have,
\begin{equation}\label{4}
\|T\|_{J}^2 = \|T^*\|_{J}^2 = \sup\limits_{k \neq 0} \frac{\|T^*k\|_{J}^2} {\|k\|_{J}^2} \leq \frac { (\beta_1 + \beta_2) \|k\|_{J}^2}{\|k\|_{J}^2} = (\beta_1 + \beta_2).
\end{equation}
Thus using equations (\ref{3}) and (\ref{4}), for every $k \in \mathcal K$ and every $\sigma \subset \mathbb N$ we obtained,
\begin{equation}
\frac{\alpha^2}{(\beta_1 + \beta_2)} \|k\|_{J}^2 \leq \sum\limits_{n \in \sigma} |[k, z_n]_{J}|^2 + \sum\limits_{n \in \sigma^c} |[k, z'_n]_{J}|^2 .
\end{equation}
The upper weaving frame bound will be achieved automatically under the given assumption. Therefore, applying theorem \ref{char}, $\{z_n\}_{n=1}^\infty$ and $\{z'_n\}_{n=1}^\infty$ are weaving frames for $\mathcal K$.
\end{proof}

The validity of the above theorem is illustrated through the following example.

\begin{ex}
Let us consider the frames 
$$
\{z_n\}_{n=1}^\infty=\{e_1, e_2, e_1, e_2, e_3, e_3, \cdots\}
$$ and 
$$
\{z'_n\}_{n=1}^\infty=\{e_1, e_1, e_2, e_2, e_3, e_3, \cdots\}
$$ for $\ell^2=\mathcal K$. Then for every $\sigma \subset \mathbb{N}$, it is straightforward to verify that the necessary and sufficient conditions of the theorem are satisfied.
\end{ex}


The following result establishes the wovenness of a Krein space frame and its image under a bounded linear operator, subject to certain conditions.
 
 \begin{prop} \label{2.8}
Let us consider a frame $\{z_n\}_{n=1}^\infty$  for Krein sapce $(\mathcal K, [.,.])$ with bounds  $\alpha$,  $\beta$ and $U$ be a bounded operator. If  $\|I - U\|_{J}^2 < \frac{\alpha}{\beta}$ then $\{z_n\}_{n=1}^\infty$ and $\{Uz_n\}_{n=1}^\infty$ are weaving frames.
\end{prop}
 
 \begin{proof}
 Clearly $U$ is invertiable and $\{Uz_n\}_{n=1}^\infty$ is a frame for $\mathcal K$.
 For every $k \in \mathcal K$ we have,      
 \begin{align*}
 \left(\sum\limits_{n \in \sigma} |[k, z_n]|^2 + \sum\limits_{n \in \sigma^c} |[k, Uz_n]|^2 \right)^\frac{1}{2} \\
 &\hspace*{-3cm} = \left(\sum\limits_{n\in\sigma} |[k, z_n]|^2 + \sum\limits_{n \in \sigma^c} |[k, z_n] - [(I - U^*)k, z_n]|^2 \right)^\frac{1}{2}\\
 & \hspace*{-3cm} \geq  \left(\sum\limits_{n \in \mathbb{N}} |[k, z_n] - \chi_{\sigma^c}(n)[(I - U^*)k, z_n]|^2 \right)^\frac{1}{2}\\
  & \hspace*{-3cm} \geq \left (\sum\limits_{n \in \mathbb{N}} |[k, z_n]|^2\right)^\frac{1}{2} - \left(\sum\limits_{n \in \mathbb{N}}|\chi_{\sigma^c}(n)[(I - U^*)k, z_n]|^2 \right)^\frac{1}{2}\\
  & \hspace*{-3cm} = \left (\sum\limits_{n \in \mathbb{N}} |[k, z_n]|^2\right)^\frac{1}{2} - \left(\sum\limits_{n \in \sigma^c}(|[(I - U^*)k, z_n]|^2 \right)^\frac{1}{2}\\
  & \hspace*{-3cm} \geq \sqrt{\alpha} \|k\|_{J} - \sqrt{\beta} \|(I - U^*)k\|_{J}\\
  & \hspace*{-3cm} \geq  (\sqrt{\alpha} - \sqrt{\beta} \|(I - U^*)\|_{J})\|k\|_{J}.
 \end{align*}
 

 %Thus $\{z_n\}_{n=1}^\infty$ and  $\{Uz_n\}_{n=1}^\infty$  forms a frame having
 
  %$$(\sqrt{\alpha} -  \sqrt{\beta} \|I - U^*\|_{J})^2 > 0$$
 %as its lower frame bound.

Again for every $k \in \mathcal K$ we have, 
 \begin{eqnarray*}
 \sum\limits_{n \in \sigma} |[k, z_n]|^2 + \sum\limits_{n \in \sigma^c} |[k, Uz'_n]|^2 
 &\leq&  \sum\limits_{n \in \mathbb{N} } |[k, z_n]|^2 + \sum\limits_{n \in \mathbb{N}} |[k, Uz'_n]|^2
 \\ &\leq& ( \beta + \|U\|^2) \|k\|^2.
 \end{eqnarray*} 
Thus $\{z_n\}_{n=1}^\infty$ and $\{Uz_n\}_{n=1}^\infty$ are weaving frames.
\end{proof}


\begin{ex}
	

  Let $k_1 = (1,0),$ and $k_2 = (0, 2) $ be two elements in the Krein sapce $\mathcal{K}=(\mathbb{R}^{2}, [., . ])$. Let us define the indefinite inner product as  $[x, y] = [(x_1, x_2), (y_1, y_2)] = x_1y_1 - x_2y_2$. The set $\{k_i\}^2_{i = 1}$  forms a frame in the krein space with bounds ($\alpha$ = 1 )and ($\beta$ = 2).
				
  Suppose $U$ is a bounded linear operator given by $U(x_1, x_2) = (x_1, 0.8x_2)$ and compute $\|I - U\|_{J} = 0.2$. Since $\|I - U\|^2_{J} = 0.04 < (\dfrac{\alpha}{\beta} = \dfrac{1}{2})$, then the Proposition \ref{2.8} ensures that $\{k_i\}^2_{i = 1}$ and $\{Uk_i\}^2_{i =1}$  are weaving frames.
	\end{ex}

\begin{ex}
	
Here we discuss a counter example of the Proposition \ref{2.8}.
	 Let us consider two elements $k_1 = (1,0),$ and $k_2 = (0, 3) $ in the Krein sapce $\mathcal{K}=(\mathbb{R}^{2}, [., . ])$ and the corresponding indefinite inner product is defined by $[x, y] = [(x_1, x_2), (y_1, y_2)] = x_1y_1 - x_2y_2$. Then the set $\{k_1, k_2\}$ forms a frame for the Krein space with bounds ($\alpha$ = 1 )and ($\beta$ = 3).
	
  Let $U$ be a bounded and  linear operator given by $U(x_1, x_2) = (-x_1, x_2)$, clearly $\{Uk_i\}^2_{i =1}$  is a frame. It is easy to verify that  $\{k_i\}^2_{i = 1}$ and $\{Uk_i\}^2_{i =1}$  are weaving frames.
However $\|I - U\|^2_{J} = 4 > (\dfrac{\alpha}{\beta} = \dfrac{1}{3})$.
	
\end{ex}

The following result establishes a condition for weaving frames through their respective frame operators.

\begin{prop}
	
	Let  $K$ = $\{z_n\}^\infty_{n =1}$ and $K'$ = $\{z'_n\}^\infty_{n= 1}$ be weaving frames  with universal frame bounds $\alpha$, $\beta$ and the corresponding frame operators  $S_K$ and $S_K'$, respectively .
	If  $\|S_K\|_J \|S_K^{-1} - S^{-1}_{K'}\|_J < \dfrac{\alpha}{\beta}$ (or $\|S_K'\|_J \|S_K^{-1} - S^{-1}_{K'}\|_J < \dfrac{\alpha}{\beta}$), then $S^{-1}_K = {\{S^{-1}_kz_n}\}^\infty_{n= 1}$ and $S^{-1}_{K'} = {\{S^{-1}_{k'}z'_n}\}^\infty_{n= 1}$
	are  weaving frames.
	
\end{prop}

\begin{proof}
	Without loss of generality let us  consider  $\|S_K\|_J \|S_K^{-1} - S^{-1}_{K'}\|_J < \dfrac{\alpha}{\beta}$. 
	For every $\sigma \subset \mathbb{N}$ and every $k \in (\mathcal{K}, [., .]$) we have,
	\begin{eqnarray*}
		&& \left (\sum\limits_{n \in \sigma} |[k, S^{-1}_Kz_n]|^2 + \sum\limits_{n \in \sigma^c} |[k, S^{-1}_{K'}z'_n]|^2 \right)^\frac{1}{2} \\
		&=& \left(\sum\limits_{n\in\sigma}|[S^{-1}_Kk, z_n]|^2 + \sum\limits_{n \in \sigma^c} |[S^{-1}_{K'}k, z'_n]|^2\right)^\frac{1}{2}\\ 
		&\geq& \left(\sum\limits_{n\in\sigma}|[S^{-1}_Kk, z_n]|^2 + \sum\limits_{n \in \sigma^c} |[S^{-1}_{K'}k +(S^{-1}_{K'} - S^{-1}_K)k, z'_n]|^2\right)^\frac{1}{2}
	%	\\&\geq&  (\sum\limits_{n\in\sigma}|[S^{-1}_Kk, z_n]|^2 + \sum\limits_{n \in \sigma^c} |[S^{-1}_Kk, z'_n]|^2)^\frac{1}{2} -  \sum\limits_{n \in \sigma^c} |[(S^{-1}_{K'} - S^{-1}_K)k, z'_n]|^2)^\frac{1}{2} \\
	\\	&\geq& \sqrt{\alpha} \|S^{-1}_Kk\| - \sqrt{\beta} \|(S^{-1}_{K'} - S^{-1}_K)k\|
\end{eqnarray*}
		On the other hand we have,
	\begin{eqnarray*}
			\\&&\hspace{-1cm}\sum\limits_{n \in \sigma} |[k, S^{-1}_Kz_n]|^2 + \sum\limits_{n \in \sigma^c} |[k, S^{-1}_{K'}z'_n]|^2 \\&\leq& \sum\limits_{n \in \mathbb{N}} |[S^{-1}_Kk, z_n]|^2 + \sum\limits_{n \in \mathbb{N}} |[S^{-1}_{K'}k, z'_n]|^2 \\&\leq& \sqrt{\beta}(\|S^{-1}_K\|^2 + \|S^{-1}_{K'}\|^2)\|k\|^2.
	\end{eqnarray*}
	%\end{eqnarray*}
	This completes the proof.
%	$$(\sum\limits_{n \in \sigma} |[k, S^{-1}_Kz_n]|^2 + \sum\limits_{n \in \sigma^c} |[k, S^{-1}_{K'}z'_n]|^2)^\frac{1}{2}$$ 
%	= $$(\sum\limits_{n\in\sigma}|[S^{-1}_Kk, z_n]|^2 + \sum\limits_{n \in \sigma^c} |[S^{-1}_{K'}k, z'_n]|^2)^\frac{1}{2}$$ 
%	$\geq$ $$(\sum\limits_{n\in\sigma}|[S^{-1}_Kk, z_n]|^2 + \sum\limits_{n \in \sigma^c} |[S^{-1}_{K'}k +(S^{-1}_{K'} - S^{-1}_K)k, z'_n]|^2)^\frac{1}{2}$$ $\geq$  
%	
%	$$(\sum\limits_{n\in\sigma}|[S^{-1}_Kk, z_n]|^2 + \sum\limits_{n \in \sigma^c} |[S^{-1}_Kk, z'_n]|^2)^\frac{1}{2} -  \sum\limits_{n \in \sigma^c} |[(S^{-1}_{K'} - S^{-1}_K)k, z'_n]|^2)^\frac{1}{2} $$ $\geq$ 
%	$$ \sqrt{\alpha} \|S^{-1}_Kk\| - \sqrt{\beta} \|(S^{-1}_{K'} - S^{-1}_K)k\|.$$ and 
%	
%	$$\sum\limits_{n \in \sigma} |[k, S^{-1}_Kz_n]|^2 + \sum\limits_{n \in \sigma^c} |[k, S^{-1}_{K'}z'_n]|^2$$ $\leq$ $$ \sum\limits_{n \in \mathbb{N}} |[S^{-1}_Kk, z_n]|^2 + \sum\limits_{n \in \mathbb{N}} |[S^{-1}_{K'}k, z'_n]|^2$$ $\leq$ $$\sqrt{\beta}(\|S^{-1}_K\|^2 + \|S^{-1}_{K'}\|^2)\|k\|^2.$$
	
\end{proof}

In the following theorem, we characterize weaving frames in Krein space using the projection operator.

\begin{thm}
	
	Let $\{z_n\}_{n=1}^\infty$ and $\{z'_n\}_{n=1}^\infty$ be two frames for Krein space $\mathcal{K},[., .]$, then the following are equivalent:
	
	\begin{enumerate}
	
	\item  The frames $\{z_n\}_{n=1}^\infty$ and $\{z'_n\}_{n=1}^\infty$ are weaving frames for $\mathcal{K}$.
	
\item	 For every $\sigma\subset \mathbb{N}$, let $Z_\sigma = \overline{span}\{z_n\}_{n \in\sigma}$,  $Z_{\sigma^{c}}= \overline{span}\{z'_n\}_{n \in\sigma^{c}}$ and let $P$ be the orthogonal projection of $\mathcal{K}$ onto $Z^\perp_{\sigma}$. Then $P|Z_{\sigma^{c}}$ is the orthogonal projection onto $Z^\perp_\sigma$ and $\{Pz'_n\}_ {n \in \sigma^c}$ is a frame for $Z^\perp_{\sigma}$.
		\end{enumerate}
\end{thm}

\begin{proof}
	
 \noindent  (\underline{1 $\implies$ 2})\\
 First we prove that for every $\sigma \subset \mathbb N$, $P|Z_{\sigma^c}$ be the orthogonal projection  onto $Z^\perp_{\sigma^c}$.
 
	Let us define $ Z_{\sigma} = \overline{span} \{z_n\}_{n \in \sigma}$, 
	$ Z_{\sigma^c} = \overline{span} \{z'_n\}_{n \in \sigma^c}$ and P be the orthogonal projection of $\mathcal{K}$ onto $Z_{\sigma}^\perp$.
	Assumed condition implies that for every $\sigma \subset \mathbb N$, $(\{z_n\}_{n \in \sigma} \cup \{z'_n\}_{n \in \sigma^c})$ is a  frame for $\mathcal{K}$.
	Since $\{z'_n\}_{n \in \sigma^c}$ spans $Z_{\sigma}^\perp$ then every $k \in Z_{\sigma}^\perp$ can be written as linear combination of $\{z'_n\}_{n \in \sigma^c}$, therefore  $P|Z_{\sigma^c}$ be the orthogonal projection  onto $Z^\perp_{\sigma}$.
	
%	Now, to show that $\{Pz'_n\}_ {n \in \sigma^c}$ is a frame for $Z_{\sigma^{c}}$.

	Since $(\{z_n\}_{n \in \sigma} \cup \{z'_n\}_{n \in \sigma^c})$ is a  frame for $\mathcal{K}$, then for every  $k \in \mathcal{K}$, there are universal constants $\alpha \leq \beta$ we have,
	$$\alpha \|k\|_{J}^2 \leq \sum\limits_{n \in \sigma} |[k, z_n]|^2 + \sum\limits_{n \in \sigma^c} |[k, z'_n]|^2 \leq \beta \|k\|_{J}^2.$$
Hence	for every $k \in Z_{\sigma}^\perp$ we have,
	$$\alpha \|k\|_{J}^2 \leq  \sum\limits_{n \in \sigma^c} |[k, z'_n]|^2 \leq \beta \|k\|_{J}^2.$$
	Since $P(z'_n) = z'_n$, then $\{Pz'_n\}_{n \in \sigma^c}$ is a frame for $Z_{\sigma}^\perp.$\\
	
	\noindent  (\underline{2 $\implies$ 1})\\
	Since $\{z_n\}_{n=1}^\infty$ is a frame  for $\mathcal{K}$ and $ Z_\sigma = \overline{span}(\{z_n\}_{n\in \sigma})$) then for every $\sigma \subset\mathbb{N}$,
	$\{z_n\}_{n\in \sigma}$ is frame for $Z_\sigma.$
	Therefore, for every $\sigma\subset \mathbb{N}$ and for every $k_\sigma \in Z_\sigma$, there exist $0<\alpha_1\leq \beta_1< \infty$ so that, 
	\begin{equation}\label{12}
	\alpha_1 \|k_\sigma\|_{J}^2 \leq \sum\limits_{n \in \sigma} |[k_\sigma, z_n]|^2  \leq \beta_1 \|k_\sigma\|_{J}^2.
	\end{equation}
	Moreover, $\{Pz_n\}_{n \in \sigma^c}$ is a frame for $Z^\perp_{\sigma}.$ Therefore, for every $\sigma\subset \mathbb{N}$ and for every $k^\perp_{\sigma} \in Z^\perp_\sigma$, there exist $0<\alpha_2\leq \beta_2< \infty$ so that
	\begin{equation}\label{13}
		\alpha_2\|k_\sigma^\perp\|_{J}^2 \leq \sum\limits_{n \in \sigma^c} |[k_\sigma^\perp, Pz'_n]|^2  \leq \beta_2 \|k_\sigma^\perp\|_{J}^2.
	\end{equation}
		Furthermore, since $P|Z_{\sigma^c}$ is the orthogonal projection  onto $Z^\perp_{\sigma}$, then from the equation (\ref{13}) we have,
		\begin{equation}\label{14}
		\alpha_2\|k_\sigma^\perp\|_{J}^2 \leq \sum\limits_{n \in \sigma^c} |[k_\sigma^\perp, z'_n]|^2  \leq \beta_2 \|k_\sigma^\perp\|_{J}^2.
	\end{equation}
	Again for every $k \in \mathcal{K}$ we have, $k = k_\sigma + k^\perp_\sigma$.
	Thus applying equations (\ref{12}) and (\ref{14}), for every $\sigma\subset \mathbb{N}$ we obtain,
	
	$$\alpha^*(\|k_\sigma^\perp\|_{J}^2 +\|k_\sigma^\perp\|_{J}^2 )\leq \sum\limits_{n \in \sigma} |[k_\sigma, z'_n]|^2 +\sum\limits_{n \in \sigma^c} |[k_\sigma^\perp, z'_n]|^2  \leq \beta^* (\|k_\sigma^\perp\|_{J}^2 +\|k_\sigma^\perp\|_{J}^2),$$
	where $\alpha^* = \min (\alpha_1, \alpha_2)$	and $\beta^* = \beta_1 + \beta_2$. Consequently, for every $\sigma\subset \mathbb{N}$ and for every $k \in \mathcal{K}$ we obtain,
	
	$$\alpha^*\|k\|_{J}^2  \leq \sum\limits_{n \in \sigma} |[k_\sigma, z'_n]|^2 +\sum\limits_{n \in \sigma^c} |[k_\sigma^\perp, z'_n]|^2  \leq \beta^* \|k\|_{J}^2.$$
This completes the proof.	
\end{proof}




\begin{defn}\cite{Gi12}\label{Bessel} 
A Bessel sequence $\{z_n\}^\infty_{n = 1}$ is said to be a $J$-frame for ($\mathcal{K}, [., .]$) if $R(T_{\sigma}P^+)$ is maximal uniformly $J$-positive subspace of $\mathcal{K}$ and $R(T_{\sigma}P^-)$ is maximal uniformly $J$-negative subspace of $\mathcal{K}$.

 
	   
	  \end{defn}

\begin{rem}\label{rem}
	Let ($\mathcal{K}, [., .]$) be a Krein space with fundamental symmetry $J$. Consider $\{z_n\}^\infty_{n = 1}$ be a Bessel sequence in $\mathcal{K}$. Let 
	$$T_{\sigma}: \ell^2(\mathbb N)\rightarrow \mathcal{K},$$ be the associated synthesis operator. Define the index sets by 
	$$\sigma_+ =\{n \in\mathbb{N} : [z_n, z_n]\geq0\}\quad and\quad \sigma_- =\{n \in\mathbb{N} : [z_n, z_n]< 0\},$$ then the Hilbert space $\ell^2(\mathbb N)$ can be orthogonally decomposed as 
	$$\ell^2(\mathbb N) = \ell^2(\sigma_+) \oplus \ell^2(\sigma_-).$$
	Let $P^+$ and $P^-$ be the orthogonal of projections from $\ell^2(\mathbb N)$ onto the subspaces $\ell^2(\sigma_+)$ and $\ell^2(\sigma_-)$, respectively. If $M_{\sigma_+} = \overline{span} \{z_n\}_{n \in \sigma_+}$ and $M_{\sigma_-} = \overline{span} \{z_n\}_{n \in \sigma_-},$
	then it is easy to check that $\overline{span} \{z_n\}_{n \in \sigma_+} \subset R(T_{\sigma}P^+) \subset M_{\sigma_+},$
	$\overline{span} \{z_n\}_{n \in \sigma_-} \subset R(T_{\sigma}P^-) \subset M_{\sigma_-} \quad.$ Thus we have, 
	$$R(T_\sigma )= R(T_{\sigma}P^+) +R(T_{\sigma}P^-).$$ 
	
\end{rem}

%\begin{defn}
	
	%Let $\{z_n\}^\infty_{n = 1}$ and $\{z'_n\}^\infty_{n = 1}$ be two J-frame for Krein space ($\mathcal {K}, [. , .])$. Then they are J-weaving frame if  each subset  $\sigma_\pm \subset \mathbb{N}$ such that  ($\mathcal{R}(T_{\sigma_+}) \cup  \mathcal{R}(T_{\sigma^{c}_+})$) is a maximal uniformly J-positive subspace of ($\mathcal{K} [., .])$ and ($\mathcal{R}(T_{\sigma_-}) \cup  \mathcal{R}(T_{\sigma^{c}_-})$) is a maximal uniformly J-negative subspace of ($\mathcal{K} [., .]$).
	
	%Let ($\mathcal{K} [. ,.]$) be Krein space with fundamental symmetry J.
	
	% Given a J--frames $\{z_n\}$ in ($\mathcal{K} [. ,.]$).
	 %consider the synthesis operator $T_{\sigma}: \ell^2(\mathbb N)\rightarrow \mathcal{K}$, . 
%	If $\sigma_+ =\{n \in\mathbb{N} : [z_n, z_n]\geq0\}$ and $\sigma_- =\{n \in\mathbb{N} : [
	%z_n, z_n]< 0\}$. 
	%Consider the orthogonal decomposition of $\ell^2(\mathbb N)$ given by $\ell^2(\mathbb N) = \ell^2(\sigma_+) \oplus \ell^2(\sigma_-)$ and denoted by $P_{\pm}$ the orthogonal projection onto $\ell^2(\sigma_{\pm})$. Also, $T_{\sigma_\pm} = T_{\sigma}P_{\pm}$ and .
	% $T_{\sigma}: \ell^2(\mathbb N)\rightarrow \mathcal{K}$  is the synthesis operator of $\{z_n\}^\infty_{n= 1}$, the J-adjoint of $T_{\sigma}$ is defined by $T^+_{\sigma}$, $T_{\sigma_+}$ and $T_{\sigma_-}$ can be easily calculated, in fact if $k \in \mathcal{K}$ : $T^+_{\sigma_\pm}k$ = ${\pm}\sum\limits_{n \in \sigma}[k, \{z_n\}^\infty_{n= 1}]e_n$, and $T^+_{\sigma} = (T_{\sigma_+} + T_{\sigma_-})^+ k = T_{\sigma_+}^+ k + T^+_{\sigma_-} k$. 
	 
	
	 
%Now, given a J--frames $\{z'_n\}$ in ($\mathcal{K} [. ,.]$).
	% consider the synthesis operator $T_{\sigma^c}: \ell^2(\mathbb N)\rightarrow \mathcal{K}$,  
	 %If $\sigma^c_+ =\{n \in\mathbb{N} : [z'_n, z'_n]\geq0\}$ and $\sigma^c_- =\{n \in\mathbb{N} : [ z'_n, z'_n]< 0\}$. 
	% Consider the orthogonal decomposition of $\ell^2(\mathbb N)$ given by $\ell^2(\mathbb N) = \ell^2(\sigma^c_+) \oplus \ell^2(\sigma^c_-)$ and denoted by $P_{\pm}$ the orthogonal projection onto $\ell^2(\sigma^c_{\pm})$. %Also,  
	 % Also, $T_{\sigma^c_\pm} = T_{\sigma^c}P_{\pm}$ and 
	 %$T_{\sigma^c}: \ell^2(\mathbb N)\rightarrow \mathcal{K}$  is the synthesis operator of $\{z_n\}^\infty_{n= 1}$, the J-adjoint of $T_{\sigma^c}$ is defined by $T^+_{\sigma^c}$, $T_{\sigma^c_+}$ and $T_{\sigma^c_-}$ can be easily calculated, in fact if $k \in \mathcal{K}$ : $T^+_{\sigma^c_\pm}k$ = ${\pm}\sum\limits_{n \in \sigma^c}[k, \{z'_n\}^\infty_{n= 1}]e_n$, and $T^+_{\sigma^c} = (T_{\sigma^c_+} + T_{\sigma^c_-})^+ k = T_{\sigma^c_+}^+ k + T^+_{\sigma^c_-} k$
	
%	If  $M_{\sigma_\pm}  = \overline{span} (\{z_n\}_{n \in \sigma_{\pm}} \cup \{z'_n\}_{n \in \sigma^c_{\pm}})$ %0r $M_{\pm}  = \overline{span} \{z_n : n \in\sigma_{\pm}\} \cup \overline {span}  \{z'_n : n \in\sigma_{\pm}\}$),
	%notice that
	%$M_{\sigma_\pm}  = \overline{span} (\{z_n\}_{n \in \sigma_{\pm}} \cup \{z'_n\}_{n \in \sigma^c_{\pm}})\subset \mathcal({R}(T_{\sigma_{\pm}}) \cup  (\mathcal{R}T_{\sigma^{c}_{\pm}}))  \subset M_{\sigma\pm}$ and ($\mathcal{R}(T_{\sigma}) \cup  \mathcal{R}(T_{\sigma^c})) = \mathcal({R}(T_{\sigma_+}) \cup  \mathcal{R}(T_{\sigma^{c}_+})) + (\mathcal{R}(T_{\sigma_-}) \cup  \mathcal{R}(T_{\sigma^{c}_-})$)  .
	
%	Moreover $\{z_n\}^\infty_{n= 1}$ and $\{z'_n\}^\infty_{n= 1}$ is a J--weaving frame, as a conseqence of its maximality, ($\mathcal({R}(T_{\sigma_{\pm}}) \cup  (\mathcal{R}T_{\sigma^{c}_{\pm}})$)  is closed. So, ($\mathcal({R}(T_{\sigma_{\pm}}) \cup  (\mathcal{R}T'_{\sigma^{c}_{\pm}})$), by [11, corollary 1.5.2(JMAA)], $M_{\sigma_+} + M_{\sigma_-} = \mathcal{K}$. Then , it follows that  $\{z_n\}^\infty_{n= 1}$ and $\{z'_n\}^\infty_{n= 1}$ is a weaving frame for the associated Hilbert space ($\mathcal{K}, <., .>$) because  ($\mathcal{R}(T_{\sigma}) \cup  \mathcal{R}(T_{\sigma^c})) = \mathcal({R}(T_{\sigma_+}) \cup  \mathcal{R}(T_{\sigma^{c}_+})) + (\mathcal{R}(T_{\sigma_-}) \cup  \mathcal{R}(T_{\sigma^{c}_-})) = M_{\sigma_+} + M_{\sigma_-}$.
	
%\end{defn}

We discuss $J$-weaving frames in Krein spaces, examining their structural properties and implications.

\begin{defn}\label{J-weaving}
	Let $\{z_n\}^\infty_{n= 1}$ and $\{z'_n\}^\infty_{n= 1}$ be two $J$-frames for Krein space ($\mathcal{K}, [., .]$) then they are said to be $J$-weaving frames if for every subset $\sigma \subset{\mathbb{N}}$,  $\{z_n\}_{n \in \sigma} \cup \{z'_n\}_{n \in \sigma^c}$ is a $J$-frame for $\mathcal{K}$.
	
If $\{z_n\}^\infty_{n= 1}$ and $\{z'_n\}^\infty_{n= 1}$ are $J$-weaving frames then for every subset $\sigma \subset\mathbb{N}, \{z_n\}_{n \in \sigma}  \cup  \{z'_n\}_{n \in \sigma^c}$ is a $J$-frame, then due to its maximality the range is closed. Thus applying \cite[Corollary 1.5.2]{An79} we have, $$R(T_{\sigma}P^+) = M_{\sigma_+} \quad and \quad R(T_{\sigma}P^-) = M_{\sigma_-} ,$$ and hence
	$$R(T_{\sigma}P^+) +R(T_{\sigma}P^-) = M_{\sigma_+} + M_{\sigma_-} = \mathcal{K}.$$ It is to be noted that $M_{\sigma_+}  = \overline{span} (\{z_n\}_{n \in \sigma_+} \cup \{z'_n\}_{n \in \sigma^c_+})$ and $M_{\sigma_-}  = \overline{span} (\{z_n\}_{n \in \sigma_-} \cup \{z'_n\}_{n \in \sigma^c_-})$.
	% because here $J$-frame this type of $\{z_n\}_{n \in \sigma}  \cup  \{z'_n\}_{n \in \sigma^c}$..
\end{defn}

%\begin{re}
	%Let $\{z_n\}^\infty_{n= 1}$ and $\{z'_n\}^\infty_{n= 1}$ such that $[z_n, z_n]\geq0$ if and only if $[z'_n, z'_n]\geq0$ and $[z_n, z_n]<0$ if and only if $[z'_n, z'_n]<0$, for $n = 1, 2, 3, \cdots $
	
%\end{re}

\begin{ex}
	If we consider $$
	\{z_n\}_{n =1}^3 = \left\{ \left(\dfrac{1}{2.16}, 0, 0 \right), \left(0, 0, \dfrac{2}{2.17} \right), \left(0, \dfrac{1}{2.18}, 0 \right)\right\}
	$$ and 
$$
\{z'_n\}_{n =1}^3 = \left\{  \left(\dfrac{100}{\sqrt{2}}, 0, \dfrac{101}{\sqrt{2}}\right), \left(\dfrac{102}{\sqrt{3}}, 0, \dfrac{103}{\sqrt{2}}\right), \left(0, \dfrac{1}{2.19}, 0\right)\right\},
$$ then it is easy to verify that $\{z_n\}_{n =1}^3$ and $\{z'_n\}_{n =1}^3$ are J-weaving frames for Krein space $\mathcal{K} = (\mathbb{R}^3, [., .])$ where  $$[(x_1, x_2, x_3), (y_1, y_2, y_3)] = x_1y_1 - x_2 y_2 + x_3 y_3.$$ 

Using the associated index sets in Remark \ref{rem}, we have the following cases:
 
{\bf Case 1:} $\sigma_+ = \{1\}$, $ M_{\sigma_+} = \overline{span} (\{z_1\} \cup \{z'_2\})$ is a maximal uniformly $J$-positive subspace of $\mathcal{K}.$ 
 
 {\bf Case 2:} $\sigma_+ = \{2\}$, $ M_{\sigma_+} = \overline{span} (\{z_2\} \cup \{z'_1\})$ is a maximal uniformly $J$-positive subspace of $\mathcal{K}.$
	
{\bf Case 3:} $\sigma_+ = \{1, 2\}$,  $ M_{\sigma_+} = \overline{span} \{z_n\}_{n = 1}^2 $ is a maximal uniformly $J$-positive subspace of $\mathcal{K}.$
	 
{\bf Case 4:}	Otherwise,  $ M_{\sigma_+} = \overline{span} \{z'_n\}_{n = 1}^2 $ is a maximal uniformly $J$-positive subspace of $\mathcal{K}.$ 
	  
It is to be noted that for every subset $\sigma_+,  M_{\sigma_+} = \overline{span} (\{z_n\}_{n \in   \sigma_+} \cup \{z'_n\}_{n \in \sigma^c_+}),$ is a maximal uniformly $J$-positive subspace of $\mathcal{K}.$ 
	  
Analogously for every subset $\sigma_-,  M_{\sigma_-} = \overline{span} (\{z_n\}_{n \in   \sigma_-} \cup \{z'_n\}_{n \in \sigma^c_-}),$ is a maximal uniformly $J$-negative subspace of $\mathcal{K}.$
	
	
\end{ex}

The following is a non-example of the definition \ref{J-weaving}, illustrating that not every weaving frames in Krein space necessarily qualify as  $J$-weaving frames.

\begin{ex}

	
	%The following is a non-example of the definition \ref{J-weaving} and also every weaving frames on Krein space  need not be $J$-weaving frames on Krein space.\\
	
	Let $(\mathcal{K} = \mathbb{C}^3, [., .])$ be a Krein space and define the associated indefinite inner product as $$[(x_1, x_2, x_3), (y_1, y_2, y_3)] = x_1\overline y_1 + x_2\overline y_2 - x_3\overline y_3.$$
Let us consider two frames 
$$ 
\{z_n\}_{n =1}^3 = \left\{ \left(3, 0, \dfrac{3}{\sqrt{2}}\right), \left(0, \dfrac{2}{\sqrt{3}}, 0\right), \left(0, 0, \dfrac{1}{\sqrt{3}}\right)\right\}
$$ and 
$$
\{z'_n\}_{n =1}^3 = \left\{ \left(1, 0, 0\right), \left(0, 3, \dfrac{3}{\sqrt{2}}\right), \left(0, 0,1\right )\right\}
$$ in $\mathbb C^3$. Then it is easy to verify that they are weaving frames, however they are not $J$-weaving frames.
	
	Here  $M_{\sigma_+}  = \overline{span} (\{z_n\}_{n \in \sigma_+} \cup \{z'_n\}_{n \in \sigma^c_+})$, $M_{\sigma_-}  = \overline{span} (\{z_n\}_{n \in \sigma_-} \cup \{z'_n\}_{n \in \sigma^c_-})$ where $\sigma_+ =\{n \in \mathbb{N}  : [z_n, z_n]\geq0\}$, $\sigma_- =\{n\in\mathbb{N} : [z_n, z_n]< 0\}$,  
	$\sigma^c_+ =\{n\in\mathbb{N}  : [z'_n, z'_n]\geq0\}$ and $\sigma^c_- = \{n \in\mathbb{N}  : [ z'_n, z'_n]< 0\}$. Thus $M_{\sigma_+}$ is not uniformly $J$-positive, whenever we consider $\sigma_+ = \{1\}$. Therefore, $\{z_n\}_{n = 1}^3$ and $\{z'_n\}_{n = 1}^3$ are not $J$-weaving frames.
	
	%Moreover, $M_{\sigma_+}  = \overline{span} (\{z_1\} \cup \{z'_2\})$  if $\{1\} \in\sigma_+$, $\{2\} \in\sigma^c_+$. But $M_{\sigma_+}$ is not uniformly J-positive, because $(3, 3, \dfrac{6}{\sqrt{2}}) \in M_{\sigma_+}$ is a (non-trival) vector with 
	%$[(3, 3, \dfrac{6}{\sqrt{2}}), (3, 3, \dfrac{6}{\sqrt{2}})] = 0$
\end{ex}





\begin{prop}\label{2.18}
If $\{z_n\}^\infty_{n= 1}$ and $\{z'_n\}^\infty_{n= 1}$ are $J$-weaving frames on Krein space $\mathcal{K}$ then for every $ \sigma_{\pm} \subset{\mathbb{N}}$ and for every $k \in M_{\sigma_\pm}$ there exist constants $ \beta_{-} \leq  \alpha_- < 0 < \alpha_+ \leq \beta_{+}$ we have, 

%$\{z_n\}_{n \in \sigma_{\pm}} \cup \{z'_n\}_{n \in \sigma_{\pm}^c}$ is a frame for the Hilbert space $M_{\pm}$.

\begin{equation}\label{19}
	\alpha_{\pm} [k, k]  \leq \sum\limits_{n \in \sigma_{\pm}} |[k, z_n]|^2 + \sum\limits_{n \in \sigma_{\pm}^c} |[k, z'_n]|^2 \leq \beta_{\pm} [k, k].
\end{equation}

\end{prop}

%\begin{ex}
%	In particular, we consider the following counter example, where Proposition \ref{2.18} is conversely not hold.\\

%Let $(\mathcal{K} = \mathbb{C}^3, [., .])$ is krein space and $$[(x_1, x_2, x_3), (y_1, y_2, y_3)] = x_1\overline y_1 + x_2\overline y_2 - x_3\overline y_3.$$


%\end{ex}

\begin{proof}
Given $\{z_n\}^\infty_{n= 1}$ and $\{z'_n\}^\infty_{n= 1}$ are $J$-weaving frames, which impiles that for every $ \sigma \subset{\mathbb{N}}$, 
$\{z_n\}_{n \in \sigma} \cup \{z'_n\}_{n \in \sigma^c}$ is a $J$-frame for $\mathcal{K}$.
Thus $R(T_{\sigma}P^+) = M_{\sigma_+}$ is a maximal uniformly $J$-positive subspace for $\mathcal{K}$. Therefore, $T_{\sigma}P^+: \ell^2(\mathbb N)\rightarrow (M_{\sigma_+}, [., .])$ is an onto map. Hence for every subset $\sigma_+ \subset \mathbb{N},  \{z_n\}_{n \in \sigma_+} \cup \{z'_n\}_{n \in \sigma_+^c}$ is a frame for ($M_{\sigma_+}, [., .]$). In particular, there exist constants $0<\alpha_+ \leq \beta_{+}$ such that equation \ref{19} holds for $M_{\sigma_+}$. A similar assertion applies
for every $\sigma_- \subset \mathbb{N}.$
\end{proof}

The following result establishes a characterization theorem for $J$-weaving frames. This characterization provides deeper insights into their structural properties within Krein spaces.

\begin{thm}\label{2.20}
	Let $\{z_n\}^\infty_{n= 1}$ and $\{z'_n\}^\infty_{n= 1}$ be weaving frames for $\mathcal{K}, [., .]$. Define the index sets $\sigma_+ =\{n \in\mathbb{N} : [z_n, z_n]\geq0\}\quad and\quad \sigma_- =\{n \in\mathbb{N} : [z_n, z_n]< 0\}$. Suppose $M_{\sigma_+}  = \overline{span} (\{z_n\}_{n \in \sigma_+} \cup \{z'_n\}_{n \in \sigma^c_+})$ and  $M_{\sigma_-}  = \overline{span} (\{z_n\}_{n \in \sigma_-} \cup \{z'_n\}_{n \in \sigma^c_-})$. Then the following are equivalent:
	\begin{enumerate}
	\item $\{z_n\}^\infty_{n= 1}$ and $\{z'_n\}^\infty_{n= 1}$ are $J$-weaving frames for $\mathcal{K}, [., .]$.
	
\item If $M_{\sigma_+} \cap M^\perp_{\sigma_+} = \{0\}\quad and\quad M_{\sigma_-} \cap M^\perp_{\sigma_-} = \{0\} $, then for every $ \sigma_{\pm} \subset \mathbb{N}$ and for every $k \in M_{\sigma_{\pm}}$, there exist constants $ \beta_{-} \leq  \alpha_- < 0 < \alpha_+ \leq \beta_{+}$ we have,
	\begin{equation}	\alpha_{\pm} [k, k]  \leq \sum\limits_{n \in \sigma_{\pm}} |[k, z_n]|^2 + \sum\limits_{n \in \sigma_{\pm}^c} |[k, z'_n]|^2 \leq \beta_{\pm} [k, k].\end{equation}
	\end{enumerate}
	
\end{thm}

\begin{proof}
	 \noindent  (\underline{1 $\implies$ 2})\\
	 Suppose $\{z_n\}^\infty_{n= 1}$ and $\{z'_n\}^\infty_{n= 1}$ are $J$-weaving frames for $\mathcal({K}, [., .]$), then using the definition of $J$-weaving frames and applying Proposition \ref{2.18}, we achieve our desired result.
	
	\noindent  (\underline{2 $\implies$ 1})\\
	If  $M_{\sigma_+} \cap M^\perp_{\sigma_+} = \{0\}$ and $M_{\sigma_-} \cap M^\perp_{\sigma_-} = \{0\}$, then both $M_{\sigma_+}$ and $M_{\sigma_-}$ are $J$-non-degenerated. Then applying (Proposition 3.8. \cite{Gi12}),  for every $ \sigma_{\pm} \subset \mathbb{N}$ and for every $k \in M_{\sigma_{\pm}}$, there exist constants $ \beta_{-} \leq  \alpha_- < 0 < \alpha_+ \leq \beta_{+}$ so that 
	$$	\alpha_{\pm} [k, k]  \leq \sum\limits_{n \in \sigma_{\pm}} |[k, z_n]|^2 + \sum\limits_{n \in \sigma_{\pm}^c} |[k, z'_n]|^2 \leq \beta_{\pm} [k, k].$$
%	and $M_{\sigma_+}$ is a uniformly $J$-positive subspace of $\mathcal{K}$ and  $M_{\sigma_-}$ is a uniformly $J$-negative subspace of $\mathcal{K}$.
	Furthermore, applying (Theorem 3.9. \cite{Gi12}) $M_{\sigma_+}$ is a maximal uniformly $J$-positive subspace of $\mathcal{K}$ and  $M_{\sigma_-}$ is a maximal uniformly $J$-negative subspace of $\mathcal{K}.$ Therefore, $\{z_n\}^\infty_{n= 1}$ and $\{z'_n\}^\infty_{n= 1}$ are $J$-weaving frames for $\mathcal{K}, [., .]$.
\end{proof}



 \section{Application to Probabilistic Erasure}
 
 In this section, we explore the applications of Krein space weaving frames in the context of probabilistic erasure, focusing on their stability in data recovery scenarios. Our discussion begins by introducing the fundamental concepts of the tensor product of elements in Krein spaces, which plays a crucial role in understanding the interaction of frame elements. This foundational discussion sets the stage for analyzing how weaving frames can be applied to mitigate the effects of erasure in probabilistic settings.
 
Let $\{z_n\}_{n=1}^\infty$ and $\{z'_n\}_{n=1}^\infty$ be weaving frames for $\mathcal K$ with the corresponding frame operators $S_1, S_2$, then for every $\sigma \subset \mathbb N$, the associated error operator is given by,
\begin{equation*}
E=\sum\limits_{n \in \sigma} \delta_n z_n \otimes S_1^{-1} z_n + \sum\limits_{n \in \sigma^c} \delta_n z'_n \otimes S_2^{-1} z'_n,
\end{equation*}
where $\delta_n$ is the standard dirac delta function. 
 
This error operator enables rapid representation while significantly minimizing computational cost. Its effectiveness is largely dependent on the selection of Krein space frames utilized in the encoding and decoding process. By optimizing these frame choices, the overall efficiency of the technique can be further enhanced.

In the following result, we demonstrate how weaving frames can effectively mitigate data loss in erasure scenarios.

\begin{thm}
Let $\{z_n\}_{n=1}^m$ and $\{z'_n\}_{n=1}^m$ be uniform tight weaving frames for the Krein space $\mathbb R^n$ with $\|z_n\|_{J} = \sqrt n = \|z'_n\|_{J}$. Suppose $\epsilon>0$ is very small number for which $\epsilon^2 \geq \frac{n}{m} \text{log} (n)$. Then for every $k \in \mathbb R^n$, there is an absolute constant $M>0$ for which we have, $$E\|\hat k-k\|_{J} \leq M \epsilon \|k\|_{J}.$$
\end{thm}

\begin{proof}
For every $k \in \mathbb R^n$ and every $\sigma \subset \{1, 2, \cdots, m\}$ we have,
\begin{equation*}
\tilde{E} k = \hat k-k = \frac{1}{m} \left [\sum\limits_{n \in \sigma} \epsilon_n [k, z_n]z_n + \sum\limits_{n \in \sigma^c} \epsilon_n [k, z'_n]z'_n \right],
\end{equation*}
where 
\begin{align*}
\hat k & = \frac{2}{m}  \left [\sum\limits_{n \in \sigma} [k, z_n]z_n + \sum\limits_{n \in \sigma^c} [k, z'_n]z'_n \right]\\ 
& =\frac{2}{m}  \left [\sum\limits_{n \in \sigma} \tilde{\delta_n}^{p_n} [k, z_n]z_n + \sum\limits_{n \in \sigma^c} \tilde{\delta_n}^{p_n} [k, z'_n]z'_n \right],
\end{align*}
 with
$$
\tilde{E}=\sum\limits_{n \in \sigma} \tilde{\delta_n}^{p_n} z_n \otimes z_n + \sum\limits_{n \in \sigma^c} \tilde{\delta_n}^{p_n} z'_n \otimes z'_n 
~\mbox{ and }~
\epsilon_n = 2\tilde{\delta_n}^{p_n} - 1
$$ is an independent Bernoulli variable that takes values $1$ and $-1$, each with probability of $\frac{1}{2}$ (see \cite{Pa99}).

Thus it is sufficient to prove that, for every $\sigma \subset \{1, 2, \cdots, m\}$,
\begin{equation*}
E\left \|  \frac{1}{m} \left [\sum\limits_{n \in \sigma} \epsilon_n z_n \otimes z_n + \sum\limits_{n \in \sigma^c} \epsilon_n  z'_n \otimes z'_n \right] \right \|_{J} \leq M \epsilon.
\end{equation*}
Applying [Lemma, \cite{Ru99}] for every $\sigma \subset \{1, 2, \cdots, m\}$, there is a positive constant $M$ we have 

\begin{align}\label{app1}
E\left \| \left [\sum\limits_{n \in \sigma} \epsilon_n z_n \otimes z_n + \sum\limits_{n \in \sigma^c} \epsilon_n  z'_n \otimes z'_n \right] \right \|_{J}
& \leq  M \sqrt{n \text{log}(n)}\\
\nonumber & \hspace{0.3cm} \left \| \left [\sum\limits_{n \in \sigma} z_n \otimes z_n + \sum\limits_{n \in \sigma^c}  z'_n \otimes z'_n \right] \right \|^{\frac{1}{2}}_{J}.
\end{align}




%\begin{equation}\label{app1}
%E\left \| \left [\sum\limits_{n \in \sigma} \epsilon_n z_n \otimes z_n + \sum\limits_{n \in \sigma^c} \epsilon_n  z'_n \otimes z'_n \right] \right \|_{J}
%\leq M \sqrt{n \text{log}(n)} \left \| \left [\sum\limits_{n \in \sigma} z_n \otimes z_n + \sum\limits_{n \in \sigma^c}  z'_n \otimes z'_n \right] \right \|_{J}^{\frac{1}{2}}.
%\end{equation}
Furthermore, since $\{z_n\}_{n=1}^m$ and $\{z'_n\}_{n=1}^m$ are uniform tight weaving frames for   $\mathbb R^n$ with $\|z_n\|_{J} = \sqrt n = \|z'_n\|_{J}$, then for every $\sigma \subset \{1, 2, \cdots, m\}$ we have,
\begin{equation*}
\mathcal I_{\mathbb R^n} =  \frac{1}{m} \left [\sum\limits_{n \in \sigma}  z_n \otimes z_n + \sum\limits_{n \in \sigma^c}  z'_n \otimes z'_n \right], 
\end{equation*}
where $\mathcal I_{\mathbb R^n}$ is the identity operator on  $\mathbb R^n$.

Hence we obtain,
\begin{equation}\label{app2}
\left \| \left [\sum\limits_{n \in \sigma} z_n \otimes z_n + \sum\limits_{n \in \sigma^c}  z'_n \otimes z'_n \right] \right \|_{J} =m.
\end{equation}
Thus using the equations (\ref{app1}) and (\ref{app2}) we obtain,
\begin{equation*}
E\left \|  \frac{1}{m} \left [\sum\limits_{n \in \sigma} \epsilon_n z_n \otimes z_n + \sum\limits_{n \in \sigma^c} \epsilon_n  z'_n \otimes z'_n \right] \right \|_{J} \leq M \sqrt{\frac{n}{m}} ~\sqrt {\text{log}(n)}.
\end{equation*}
Therefore, we achieved our result due to the fact $\epsilon^2 \geq \frac{n}{m} \text{log} (n)$.
\end{proof}

 
 

\section*{Acknowledgment}
The authors acknowledge the Department of Mathematics at SRM University, {\it AP} - Andhra Pradesh, for providing an enriching academic environment to carry out the research. The second author expresses his gratitude to his advisors, Prof. Saikat Mukherjee (NIT Meghalaya, India) and Prof. Jaydeb Sarkar (ISI Bangalore, India), for their valuable suggestions regarding his research.

%\section*{Data availability}
%This manuscript is prepared with the assurance that no extraneous data has been incorporated, emphasizing the integrity and reliability of the presented information.



%\section*{Declarations}
%Our commitment to transparency and ethical conduct ensures that there is no conflict of interests.




\begin{thebibliography}{9}

\bibitem{An79} T. Ando,  \emph{Linear Operators on Krein Spaces}, Hokkaido University, Sapporo, Japan, 1979.

\bibitem{Ba18} R. Balan, K. A. Okoudjou, and A. Poria, On a Problem by Hans Feichtinger, \emph{Operators and Matrices},  \textbf{12} (2018),  881--891.

\bibitem{Fi06} P. Balazs, H. G. Feichtinger, M. Hampejs and G. Kracher, Double preconditioning for Gabor frames, \emph{IEEE Transactions on signal processing},  \textbf{54} (2006),  4597--4610.

\bibitem{Be16} T. Bemrose, P. G. Casazza, K. Gr{\"o}chenig, M.C. Lammers and R. G. Lynch, Weaving frames, Operators and Matrices, 10 (2016), 1093--1116.

\bibitem{Fi98} H. Bolcskei, F. Hlawatsch and H. G. Feichtinger, Frame-theoretic analysis of oversampled filter banks, \emph{IEEE Transactions on signal processing},  \textbf{46} (1998),  3256--3268.

\bibitem{Bo06}
M. Bownik and D. Speegle,  The Feichtinger Conjecture for Wavelet  Frames, Gabor Frames and
Frames of Translates,  \emph{Canad. J. Math.}, \textbf{58} (2006), 1121--1143.


\bibitem{christensen2003introduction} O. Christensen,  \emph{An Introduction to Frames and Riesz Bases}, Birkhauser, Boston, 2003.

\bibitem{3}%17
O. Christensen, S. Datta and R. Y. Kim, Equiangular frames and generalizations of the Welch bound to dual pairs of frames, \emph{Linear and Multilinear Algebra},  \textbf{68} (2020), 2495--2505.

\bibitem{Cz02}
W. Czaja and Z. Rzeszotnik,  Pseudodifferentia Operators and Gabor Frames: Spectral Asymptotics,  \emph{Mathematische Nachrichten}, \textbf{233} (2002), 77--88.

\bibitem{Cz06}
W. Czaja, D. Speegle and G. Kutyniok,  The geometry of sets of parameters of wave packet frames,  \emph{Appl. Comput. Harmon. Anal}, \textbf{20} (2006), 108--125.


\bibitem{2} I. Daubechies, A. Grossmann and Y. Meyer, Painless nonorthogonal expansions, \emph{J. Math. Phys.}, \textbf{27} (1986), 1271 -- 1283.

 \bibitem{1}%17
Jyoti, Deepshikha, L. K. Vashisht and G. Verma, Sums of matrix valued wave packet frames in $L^2(\mathbb R^d, \mathbb C^{s \times r})$,  \emph{Glasnik Mathematicki}, \textbf{ 53} (2018), 153--177.

\bibitem{duffin1952class} R. S. Duffin and A. C. Schaeffer, A class of nonharmonic Fourier series, \emph{Trans. Am. Math. Soc.}, \textbf{72} (1952), 341 -- 366.


%\bibitem{Ch20} O. Christensen, M. Hasannasab and F. Philipp, Frame properties of operator orbits, \emph{Math. Nachr.}, \textbf{ 293} (2020), 52--66.



%\bibitem{Nikolski} N. K. Nikolski, \emph{Rudiments de la theorie du contr$\hat{o}$le vision op${\Acute{e}}$ratorielle, Publ. ${\Acute{E}}$cole Doctorale Mat. Bordeaux, Cours Post DEA}, 1994.

%\bibitem{Nagy et al.} B. S. Nagy, C. Foias, H. Bercovici and L. Kerchy, \emph{Harmonic Analysis of Operators on Hilbert Space}, Springer, 2010.

%\bibitem{Cath_Sola} C. B\'{e}n\'{e}teau, D. Khavinson, C. Liaw, D. Seco and A. A. Sola, Orthogonal polynomials, reproducing kernels, and zeros of optimal approximants,  \emph{J. London Math. Soc.}, \textbf{94} (2016), 726--746.



%\bibitem{Paulsen_Raghupathi} V. I. Paulsen and M. Raghupathi,  \emph{ An Introduction to the theory of Reproducing Kernel Hilbert Spaces}, Cambridge University Press, UK, 2016.


%\bibitem{Wick11} N.  Arcozzi, R.  Rochberg, E. Sawyer and B. D. Wick,  Distance functions for reproducing kernel Hilbert spaces, \emph{Function spaces in modern analysis}, 25--53, 2011.

%\bibitem{Wick14} M. Mitkovski and B. D. Wick,  A reproducing kernel thesis for operators on Bergman-type function spaces, \emph{Journal of Functional Analysis}, \textbf{267} (2014), 2028--2055.

%\bibitem{Wick07} S. Petermichl, S. Treil and B. D. Wick,  Carleson potentials and the reproducing kernel thesis for embedding theorems, \emph{Illinois Journal of Mathematics}, \textbf{51} (2007), 1249--1263.

%\bibitem{benedetto1998theory} J. Benedetto, S. Li : The theory of multiresolution analysis frames and applications to filter banks, Appl. Comput. Harmon. Anal., 5(4), 389 -- 427 (1998).

%\bibitem{bakic2005parseval} D. Baki{\'c}, I. Krishtal and E. N. Wilson : Parseval frame wavelets with $E^{(2)}_n$ -dilations, Appl. Comput. Harmon. Anal., 19(3), 386 -- 431 (2005)
\bibitem{Ga} K. Esmeral Garc\'{i}a and  E. Wagner, Frames in Krein spaces, Preprint.

\bibitem{Fi97} H. G. Feichtinger and K Gröchenig, Gabor frames and time-frequency analysis of distributions, \emph{Journal of Functional Analysis},  \textbf{146} (1997),  464--495.


\bibitem{Fi00} H. G. Feichtinger and A.  Janssen, Validity of WH-frame bound conditions depends on lattice parameters, \emph{Applied and Computational Harmonic Analysis},  \textbf{8} (2000),  104--112.

\bibitem{Gi12} J. I. Giribet, A. Maestripieri, F. Martínez Per\'{i}a and P. G. Massey, On frames for Krein spaces, \emph{Journal Math. Anal. App.}, \textbf{393} (2012), 122--137.

\bibitem{Gu18}
A. Gumber and N. K. Shukla, Orthogonality of a pair of frames over locally compact abelian groups,
\emph{Journal of Mathematical Analysis and Applications}, \textbf{458} (2018), 1344--1360.

\bibitem{Ha07} D. Han, K. Kornelson, D. Larson, E. Weber, Frames for undergraduates, In Student mathematical library, 2007. %https://doi.org/10.1090/stml/040

\bibitem{Ha00} D. Han, D.R. Larson, Frames, bases and group representations, \emph{Mem. Am. Math. Soc.} \textbf{147}(697)(2000). %https://doi.org/10.1090/memo/0697

\bibitem{He04}  C. Heil, P. E. T. Jorgensen, D. R . Larson,  Wavelets, frames and operator theory. In Contemporary mathematics - American Mathematical Society (2004). %https://doi.org/10.1090/conm/345

\bibitem{Kaf09}%17
V. Kaftal, D. R. Larson and S. Zhang, Operator-Valued Frames,  \emph{Transaction of AMS.}, \textbf{361} (2009), 6349--6385.



 %\bibitem{2}%17D. Li, J. Leng: Operator representations of $g$-frames in Hilbert spaces, Linear and Multilinear Algebra. 68(9), 1861--1877 (2020).
 
 
 \bibitem{Na12} T. C. E. Nambudiri and K. Parthasarathy, Generalised Weyl–Heisenberg frame operators, \emph{Bull. Sci. Math.}, \textbf{136} (2012), 44--53.
 
 \bibitem{Na13} T. C. E. Nambudiri and K. Parthasarathy, A characterisation of Weyl–Heisenberg frame operators, \emph{Bull. Sci. Math.}, \textbf{137} (2013), 322--324.
 
 \bibitem{Pa99} E. Parzen, Stochastic Processes, \emph{SIAM}, \textbf{24} (1999).

\bibitem{Po20} A. Poria, Positive Weight Function and Classification of $g$-Frames, \emph{Numerical Functional Analysis and Optimization},  \textbf{41} (2020), 950--968.













% \bibitem{2}%17
%D. Li, J. Leng, Operator representations of $g$-frames in Hilbert spaces,  \emph{Linear and Multilinear Algebra}, \textbf{ 68} (2020), 1861--1877.

%\bibitem{3}%17
%O. Christensen, S. Datta, R. Y. Kim,  Equiangular frames and generalizations of the Welch bound to dual pairs of frames,  \emph{Linear and Multilinear Algebra}, \textbf{68} (2020), 2495--2505.

%\bibitem{4}%17
%X. Guo, Y. Chen,  On some operators and dilations of frame generator and dual pair of frame generators of two structured unitary systems,  \emph{Linear and Multilinear Algebra}, \textbf{70} (2022), 830--842.

%\bibitem{Be03} A. Ben and N. E. Thomas,  \emph{Generalized Inverses, Theory and Applications}, Springer, New York, 2003.
	

%\bibitem{Do66}%17
%R. G. Douglas, On majorization, factorization and range inclusion of operators on {H}ilbert space,  \emph{Proc. Amer Math. Society}, \textbf{17} (1966), 413--415.






\bibitem{Ru99} M. Rudelson, Random Vectors in the Isotropic Position, \emph{Journal Func. Anal.}, \textbf{164} (1999), 60--72.



%\bibitem{Frican-2001}
%E. Fricain,  Bases of reproducing kernel in model space,   \emph{J. Operator Theory}  \textbf{46} (2001), 517--543.

%\bibitem{Bauer} W. Bauer, R. Duduchava, S. Grudsky and M. A. Kaashoek,   \emph{Operator Algebras, Toeplitz Operators and Related Topics}, Birkhauser, 2020.

%\bibitem{Rosenthal} R. Avenda$\tilde{n}$o and P.  Rosenthal,  \emph{ An Introduction to Operators on the Hardy-Hilbert Space. Graduate Texts in Mathematics}, Springer, New York, 2007.

%\bibitem{Cab23}
%C. Cabrelli, U. Molter and D. Suarez,  Frames of iterations and vector-valued model spaces, preprint (2023),  arXiv:2203.01301v3[math.FA].
 
%\bibitem{Kilic95}
%S. Kilic:  The Berezin Symbol and Multipliers of Functional Hilbert Spaces,  Proc. AMS.  123, 3687--3691 (1995).
 
%\bibitem{Kavarev02}
%M. T. Kavarev:  Berezin symbols and Schatten Von-Neumann Classes,  Math. Notes.  72, 185--192 (2002).
 
%\bibitem{Clark72}
%D. Clark,  One dimensional perturbations of restricted shifts,   \emph{J. Math. Anal. Appl.},  \textbf{25} (1972), 169--191.
 
%\bibitem{Debnath98}
%L. Debnath and P. Mikusinki,   \emph{Introduction to Hilbert spaces with applications},  Academic Press 1998.

%\bibitem{Atteia99}
%M. Atteia and J. Gaches,  \emph{ Approximation hilbertienne: Splines, Ondelettes, Fractales},  Presses Universitaires de Grenoble, 1999.

%\bibitem{Tau49}
%O. Taussky,  A remark concerning the characteristic roots of the finite segments of the Hilbert matrix,  \emph{Quart. J. Math. Oxford Ser.}, \textbf{20} (1949), 80--83.

%\bibitem{Kato57}
%T. Kato,  On the Hilbert matrix,  \emph{Proc. Amer. Math. Soc.}, \textbf{8} (1957), 73--81.

%\bibitem{Rosen58}
%M. Rosenblum,  On the Hilbert matrix II,  \emph{Proc. Amer. Math. Soc.}, \textbf{9} (1958), 581--585.





\end{thebibliography}
\end{document}
